\documentclass[12pt,a4paper]{article}
\usepackage[T1]{fontenc}
\usepackage[ngerman]{babel}
\usepackage{lmodern}
\usepackage{csquotes}
\usepackage{amsmath,amssymb}
\usepackage{geometry}
\geometry{a4paper, margin=2.5cm}
\usepackage{parskip}
\usepackage{xcolor}
\usepackage[colorlinks=true, linkcolor=blue!60!black, urlcolor=blue!60!black, citecolor=blue!60!black]{hyperref}
\usepackage[backend=biber, style=numeric]{biblatex}
\usepackage[toc, acronym, nonumberlist, style=long]{glossaries}
\makeglossaries
\sloppy

\begin{filecontents*}{\jobname.bib}
@misc{bsi2023socialengineering,
  author = {{Bundesamt für Sicherheit in der Informationstechnik}},
  title = {Social Engineering: Die Schwachstelle Mensch im Fokus},
  year = {2023}
}
\end{filecontents*}

\addbibresource{\jobname.bib}

\newglossaryentry{socialengineering}{
  name=Social Engineering,
  description={Die gezielte Manipulation von Personen durch psychologische Tricks, um an vertrauliche Informationen zu gelangen oder Handlungen zu erzwingen.}
}

\newglossaryentry{pretexting}{
  name=Pretexting,
  description={Eine Social-Engineering-Methode, bei der eine erfundene Geschichte oder Identität genutzt wird, um das Vertrauen des Opfers zu gewinnen.}
}

\begin{document}

% --- TITELBLATT ---
\begin{titlepage}
    \centering
    \vspace*{2cm}
    
    {\Huge \textbf{Skript: Social Engineering} \par}
    \vspace{1.5cm}
    {\Large \textbf{Kantonsschule im Lee} \par}
    \vspace{2cm}
    
    \rule{\linewidth}{0.5mm}
    \vspace{2cm}
    
    {\large \textbf{Thema:} Social Engineering, Angriffsvektoren \& Abwehrmassnahmen \par}
    \vspace{0.5cm}
    {\large \textbf{Fach:} Informatik \par}
    \vspace{0.5cm}
    {\large \textbf{Autor:} Laurin Ott \par}
    
    \vfill
    
    {\large \textbf{Datum:} \today \par}
\end{titlepage}

% --- INHALTSVERZEICHNIS ---
\newpage
\tableofcontents
\newpage

\section{Einleitung}
In der IT-Sicherheit gilt der Mensch oft als das schwächste Glied in der Sicherheitskette. Während technische Schutzmassnahmen wie Firewalls und Verschlüsselung hochentwickelt sind, nutzen Angreifer beim \gls{socialengineering} gezielt menschliche Eigenschaften wie Hilfsbereitschaft, Respekt vor Autoritäten oder Naivität aus, um Sicherheitsbarrieren zu umgehen \cite{bsi2023socialengineering}.

\section{Grundlegende Angriffsformen}
\gls{socialengineering} umfasst eine Vielzahl von psychologischen Angriffsmethoden, die im Folgenden näher betrachtet werden:
\begin{itemize}
    \item \textbf{Pretexting:} Der Angreifer erfindet eine glaubwürdige Geschichte oder Identität (zum Beispiel als IT-Support oder Vorgesetzter), um an vertrauliche Daten zu gelangen.
    \item \textbf{Tailgating/Piggybacking:} Ein unbefugter Dritter verschafft sich physischen Zugang zu einem gesicherten Gebäude, indem er einer berechtigten Person durch eine Tür folgt.
    \item \textbf{Baiting:} Dem Opfer wird ein physisches Medium (wie ein präparierter USB-Stick) oder ein digitaler Köder ausgelegt, dessen Neugier geweckt werden soll.
    \item \textbf{Quid Pro Quo:} Dem Opfer wird eine Gegenleistung (zum Beispiel technische Hilfe) für die Preisgabe von Zugangsdaten oder Informationen versprochen.
\end{itemize}
                                                   
\section{Psychologische Mechanismen}
Angreifer setzen gezielt psychologische Hebel ein, um rationale Entscheidungsprozesse zu umgehen:
\begin{itemize}
    \item \textbf{Autorität:} Das Ausnutzen von Respekt vor Vorgesetzten oder offiziellen Stellen.
    \item \textbf{Dringlichkeit:} Erzeugen von Zeitdruck, damit das Opfer nicht lange nachdenkt.
    \item \textbf{Hilfsbereitschaft:} Appell an das Gute im Menschen, einem Kollegen oder einer hilfssuchenden Person schnell zu assistieren.
\end{itemize}

\section{Best Practices und Prävention}
Um sich und Unternehmen vor \gls{socialengineering} zu schützen, sind folgende Massnahmen zentral:

\begin{enumerate}
    \item \textbf{Sensibilisierung und Schulung:} Regelmässige Aufklärung der Mitarbeitenden über gängige Betrugsmaschen.
    \item \textbf{Kritische Prüfung:} Identitäten von Personen (insbesondere bei telefonischen oder digitalen Anfragen) stets über offizielle Kanäle verifizieren.
    \item \textbf{Klare Prozesse:} Eindeutige Richtlinien für die Vergabe von Zugangsdaten und physischen Zutritten etablieren.
\end{enumerate}

\printbibliography

\printglossary[type=\acronymtype, title=Akronyme]
\printglossary[title=Glossar]

\end{document}